\DocumentMetadata{
  pdfversion=2.0,
  pdfstandard=ua-2,
  lang=en-US,
  tagging=on,
  tagging-setup={math/setup=mathml-SE,tabsorder=structure}
}
\documentclass[11pt,letterpaper]{article}
\usepackage[margin=0.68in,includefoot]{geometry}
\usepackage{newcomputermodern}
\usepackage{amsmath,amscd,mathtools}
\usepackage{tikz,adjustbox}
\providecommand{\square}{\mdlgwhtsquare}
\usepackage[hidelinks]{hyperref}
\hypersetup{
  pdftitle={Numerical Analysis PhD Exam, January 2009},
  pdfauthor={University of Florida Department of Mathematics},
  pdfsubject={Graduate mathematics examination},
  pdfdisplaydoctitle=true,
  bookmarksopen=true
}
\setcounter{secnumdepth}{0}
\setlength{\parindent}{0pt}
\setlength{\parskip}{0.28em}
\setlength{\emergencystretch}{3em}
\setlength{\leftmargini}{2.15em}
\setlength{\leftmarginii}{2.65em}
\setlength{\leftmarginiii}{3em}
\pagestyle{empty}
\raggedbottom
\allowdisplaybreaks
\NewTaggingSocketPlug{block/list/label}{examproblem}{%
  \tagstructbegin{tag=Lbl}\tagstructend%
  \tagstructbegin{tag=\LBody}%
  \tagstructbegin{tag=\ProblemHeadingTag,title-o={\ProblemHeadingText}}%
  \tagmcbegin{tag=\ProblemHeadingTag,actualtext={\ProblemHeadingText}}%
  #2%
  \tagmcend%
  \tagstructend%
}
\newenvironment{examproblems}[2]{%
  \def\ProblemHeadingTag{#2}%
  \AssignTaggingSocketPlug{block/list/label}{examproblem}%
  \begin{enumerate}%
  \setcounter{enumi}{#1}%
  \renewcommand{\labelenumi}{\textbf{\theenumi.}}%
  \setlength{\topsep}{0.35em}%
  \setlength{\partopsep}{0pt}%
  \setlength{\itemsep}{0.48em}%
  \setlength{\parsep}{0.18em}%
}{\end{enumerate}}
\newenvironment{examsubparts}{%
  \AssignTaggingSocketPlug{block/list/label}{default}%
  \begin{enumerate}%
  \setlength{\topsep}{0.18em}%
  \setlength{\partopsep}{0pt}%
  \setlength{\itemsep}{0.16em}%
  \setlength{\parsep}{0.1em}%
}{\end{enumerate}}
\newenvironment{examinstructions}{%
  \par\begingroup\itshape
}{%
  \par\endgroup\vspace{0.18em}
}
\makeatletter
\renewcommand\subsection{\@startsection{subsection}{2}{0pt}%
  {-1.25ex plus -.25ex minus -.1ex}%
  {0.55ex plus .1ex}%
  {\normalfont\large\bfseries}}
\renewcommand{\@maketitle}{%
  \newpage\null\vskip -1em%
  {\footnotesize\bfseries
    \href{https://www.ufl.edu/}{University of Florida}, \href{https://math.ufl.edu/}{Department of Mathematics}\par}%
  \vskip -0.5em%
  \tagstructbegin{tag=H1,title={Numerical Analysis PhD Exam, January 2009}}%
  \tagpdfparaOff%
  \tagmcbegin{tag=H1}%
    {\LARGE\bfseries\href{https://gma.math.ufl.edu/exams/numerical-analysis-phd/}{Numerical Analysis PhD Exam}, January 2009\par}%
  \tagmcend%
  \tagpdfparaOn%
  \tagstructend%
  \vskip 0.25em%
  \tagmcbegin{artifact}%
    \hrule height 1pt%
  \tagmcend%
  \vskip 1em%
}
\makeatother
\title{Numerical Analysis PhD Exam, January 2009}
\author{}
\date{}
\begin{document}
\maketitle
\thispagestyle{empty}
\begin{examinstructions}
Answer any 8 questions.
\end{examinstructions}
\begin{examproblems}{0}{H2}
\def\ProblemHeadingText{Problem 1.}
\item
Let \(M=\begin{bmatrix}A & B^t\\ B & C\end{bmatrix}\) be real symmetric and positive definite matrix. Let \(S=C-BA^{-1}B^t\) be the Schur complement. Prove that~\(S\) is symmetric and positive definite.
\def\ProblemHeadingText{Problem 2.}
\item
Let \(\lVert\cdot\rVert_p\) for \(1\le p\le\infty\) denote the norm on \(m\times n\) matrices induced by the \(\ell^p\)-norm on vectors in \(\mathbb{C}^m\) and \(\mathbb{C}^n\). Prove that
\[
\lVert A\rVert_2\le \sqrt{\lVert A\rVert_1\lVert A\rVert_\infty},\qquad\text{for all }A\in\mathbb{C}^{m\times n}.
\]
\def\ProblemHeadingText{Problem 3.}
\item
Show that if~\(P\) and~\(Q\) be Hermitian positive definite matrices satisfying \(x^*Px\le x^*Qx\), for all \(x\in\mathbb{C}^n\), then \(\lVert P\rVert_F\le\lVert Q\rVert_F\).
\def\ProblemHeadingText{Problem 4.}
\item
Let~\(T\) be any square matrix and let \(\lVert\cdot\rVert\) denote any induced norm. Prove that \(\displaystyle\lim_{n\to\infty}\lVert T^n\rVert^{1/n}\) exists and equals \(\displaystyle\inf_{n=1,2,\ldots}\lVert T^n\rVert^{1/n}\).
\def\ProblemHeadingText{Problem 5.}
\item
Let~\(P\) and~\(Q\) be two \(m\times m\) orthogonal projectors. Prove that \(\lVert P-Q\rVert_2\le 1\).
\def\ProblemHeadingText{Problem 6.}
\item
Let \(x_0,x_1,\ldots,x_n\) be distinct points in a finite interval \([a,b]\) and \(f\in C^1[a,b]\). Show that for any given \(\varepsilon>0\) there exists a polynomial~\(p\) such that \(\lVert f-p\rVert_\infty<\varepsilon\) and \(p(x_i)=f(x_i)\) for all \(i=0,1,\ldots,n\) (where \(\lVert\cdot\rVert_\infty\) denotes the \(L^\infty(a,b)\)-norm).
\def\ProblemHeadingText{Problem 7.}
\item
Let \(\hat f(s)\) be the continuous Fourier transform of \(f(t)\in L^2[-\infty,\infty]\). Suppose further that \(\hat f(s)=0\) for \(|s|>\pi\). Derive the interpolation formula
\[
f(t)=\sum_{k=-\infty}^{\infty}f(k)\frac{\sin(\pi(t-k))}{\pi(t-k)}.
\]
\def\ProblemHeadingText{Problem 8.}
\item
Let \(L_n(f)\) be the Lagrange polynomial interpolating a function~\(f\) at nodes \(a=x_0<x_1<\cdots<x_n=b\).
\begin{examsubparts}
\item[\textnormal{(a)}]
Give the formula for \(L_n(f)\).
\item[\textnormal{(b)}]
Prove that there is a unique polynomial of degree at most~\(n\) interpolating~\(f\) at the nodes.
\item[\textnormal{(c)}]
State and prove the formula for the error \(f(x)-L_n(f)(x)\) when \(f(x)\in C^{n+1}[a,b]\).
\end{examsubparts}
\def\ProblemHeadingText{Problem 9.}
\item
Assume that~\(g\) is continuously differentiable real-valued function and \(a\le g(x)\le b\) on \([a,b]\). Show the following:
\begin{examsubparts}
\item[\textnormal{(a)}]
There is an \(\alpha\in[a,b]\) such that \(g(\alpha)=\alpha\).
\item[\textnormal{(b)}]
If \(|g'(x)|<1\) on \([a,b]\), then there is only one fixed point on the interval \([a,b]\).
\item[\textnormal{(c)}]
If we generate iterates using the recurrence \(x_{n+1}=g(x_n)\) starting from some \(x_0\in[a,b]\), then we have
\[
|\alpha-x_n|\le \frac{\lambda^n}{1-\lambda}|x_1-x_0|\qquad\text{where}\qquad \lambda=\max_{x\in[a,b]}|g'(x)|.
\]
\end{examsubparts}
\def\ProblemHeadingText{Problem 10.}
\item
Let \(\{\phi_n(x):n\ge 0\}\) be an orthogonal family of polynomials on the interval \((a,b)\) with weight function \(w(x)\ge 0\) for all \(x\in[a,b]\). Show that the polynomial \(\phi_n(x)\) has exactly~\(n\) distinct real roots in the open interval \((a,b)\).
\end{examproblems}
\end{document}
